\subsectiontitle{Markov Decision Process}

In this section, we move away from bandit models towards a setting with a linear Markov transition model, which includes a state variable. 

\subsubsectiontitle{True Environment}

We now define a state variable: $s_{i,t} \in \mathbb{R}$. (We also remove context variable $c_i$ that was used in the contextual bandit.) The state variable can be thought of as anything that can be measured on a continuous scale from $-\infty$ to $\infty$—for example, happiness. 

initial state (\enquote{happiness}): $s_{i,1} \sim \mathcal{N}(\mu_{s_0} \sigma_{s_0}^2)$ everyone's initial happiness is a random draw from this distribution

state transitions: $s_{i,t+1} \mid s_{i,t}, a_{i,t}
\sim
\mathcal{N}
\left(
\alpha_i + \beta_i s_{i,t} + \gamma_i a_{i,t}
\,
\sigma_s^2
\right).$

% kellman filtering / bayesian linear regression

% sigma should be a function of the s
% design matrix needs to be full rank, if not full rank need to do ridge regression
% intercept always 1 c_i is always 1
% talk to gemini 

same for all $i$

$\alpha_{i}$: starting happiness

$\beta_{i}$: contribution of current happiness to next time pont happiness

$\gamma_{i}$: contribution of state to next time point happiness

$\sigma_{s}$: random happiness fluctuations

$\alpha_i, \beta_i, \gamma_i, \sigma^2_s, \sigma^2_r$ are \textit{constant} across the entire population.

rewards are determined as follows: 
$r_{i,t}
=
\theta_i^\top
\begin{pmatrix}
1\\
s_{i,t}\\
a_{i,t}\\
s_{i,t}a_{i,t}
\end{pmatrix}
+
\eta_{i,t}, $ $\eta_{i,t} \sim \mathcal{N}(0,\sigma_r^2)$ (random reward fluctuations)

$\theta_i =
\begin{pmatrix}
\theta_{i,0} \\
\theta_{i,1} \\
\theta_{i,2} \\
\theta_{i,3}
\end{pmatrix}$, $\theta_i \sim \mathcal N(\mu_\theta,\Sigma_\theta)$ varies by person and drawn from underlying population distribution

\subsubsectiontitle{Action Selection Strategy}

Produces four estimates for each participant at each time point: $\hat{\theta}_{i,0,t},\; \hat{\theta}_{i,1,t},\; \hat{\theta}_{i,2,t},\; \hat{\theta}_{i,3,t}$ and four measures of uncertainty: $\hat{\Sigma}_{\theta,i,t}
=
\operatorname{diag}
\left(
\hat{\sigma}^2_{\theta_{i,0,t}},
\hat{\sigma}^2_{\theta_{i,1,t}},
\hat{\sigma}^2_{\theta_{i,2,t}},
\hat{\sigma}^2_{\theta_{i,3,t}}
\right).$

\subsubsectiontitle{Decision-Making Algorithm}

Decision-making algorithm receives eight numbers (four estimated means, four measures of uncertainty) per participant per time point.

The action selection strategy believes...

$\hat{\theta}_{i,t} \sim
\mathcal{N}(\theta_i,\Sigma_{\theta,i,t})$

$\theta_i \sim
\mathcal{N}(\mu_\theta,\Sigma_\theta)$

$\Sigma_\theta =
\operatorname{diag}
(\sigma^2_{\theta,0},\sigma^2_{\theta,1},\sigma^2_{\theta,2},\sigma^2_{\theta,3})$